% =====================================================================
% APPENDIX A (MLLB): use of AI code assistants
% Source: docs/AI_USAGE_LOG.md (Bounty candidates B4, B5, B8, B9, B18; one report per error).
% =====================================================================
\appendixbanner{APPENDIX A: USE OF AI CODE ASSISTANTS}
\addcontentsline{toc}{section}{Appendix A: Use of AI code assistants}\label{app:ai-code}

\subsection*{Workflow}

The environment and the control rules are in Section~\ref{sec:ai-env}. For each block the assistant
proposed the decisions, we accepted or changed them, the assistant wrote the code and its tests and ran
them on a CPU copy with synthetic data, and we ran the tests, the smoke run, the real runs and the report
on the laptop, pasting the output back. After the runs the assistant recomputed the main metrics from
the saved predictions and compared them with the summaries before writing the block note. Every step is
logged in \texttt{docs/AI\_USAGE\_LOG.md}.

\subsection*{Errors detected (candidates for the AI Error Bounty)}

Every assistant error that we corrected is an entry of \texttt{docs/AI\_USAGE\_LOG.md} (B1--B19), and each
one is documented in exactly one of our four reports. The five sheets below are the code errors of my
blocks and of the shared work that I detected or that the assistant found in my runs:
\begin{enumerate}[leftmargin=1.4em,itemsep=1pt]
  \item \textbf{Out of GPU memory when generating with a long context} (B8; detected by me; Sheet C1).
  \item \textbf{A figure legend placed over a data point} (B18; detected by me; Sheet C2).
  \item \textbf{Commands given without changing to the repository folder} (B9; detected by me; Sheet C3).
  \item \textbf{A wrong cause given for the ``uncommitted code'' flag} (B5; found from my terminal output;
        Sheet C4).
  \item \textbf{The smoke mode of the Dense report stopped} (B4; found from the missing outputs of my smoke
        run; Sheet C5).
\end{enumerate}
The other code errors of the project (tensors on two devices in the zero-shot code, stale copies of files
written to the folder, a wrong code comment in the Transformer and a BibTeX error missed at compile time)
are documented in Mar\'ia's MLLB report.

\subsection*{AI Error Bounty sheets}

\subsubsection*{Sheet C1. Generation ran out of GPU memory (Claude)}
\begin{description}[style=nextline,leftmargin=1.2em]
  \item[Prompt] Write the complete generative block: training, generation with several sampling settings,
        statistics and figures.
  \item[Faulty response] \texttt{generate} used batches of 64 sequences for every context length. With
        a context of 500 tokens, the attention matrices of 64 windows (64 $\times$ 4 heads $\times$ 500
        $\times$ 500 floats, several per layer) do not fit in 4\,GB. Training memory had been measured with
        \texttt{estimate}, generation memory had not, and the synthetic tests ran on the CPU.
  \item[Detected by] Me, when the report step of my generative block failed on the laptop.
  \item[Proof] I ran \texttt{python -m src.generative\_experiments report}; it stopped with
        \texttt{CUDA error: out of memory} in \texttt{scores = q @ k.transpose(-2, -1)} while generating
        with the 5\,s context.
  \item[Correction] Generation batches of $64\,(250/L)^2$ windows for $L > 250$ and the GPU freed between
        models; a test was added and the report was run again. The trained models were not affected.
  \item[Reproduction] Remove the line that shrinks \texttt{batch} in \texttt{src/generative.py::generate}
        and run \texttt{report} on the laptop: it stops at the 5\,s context.
  \item[Reflection] The assistant knew that attention memory grows with $L^2$ and still checked only the
        training path. A test on the CPU cannot reveal a memory limit of the GPU.
\end{description}

\subsubsection*{Sheet C2. A legend over a data point (Claude)}
\begin{description}[style=nextline,leftmargin=1.2em]
  \item[Prompt] Write the final comparison (task T15): table and overview figure of every final model.
  \item[Faulty response] \texttt{src/comparison.py::figure\_overview} fixed the legend at
        \texttt{loc="lower left"}, the corner where the lowest row, zero-shot CLIP (macro AUROC near 0.5),
        is drawn. The assistant inspected the figure on synthetic and on real data without noticing it.
  \item[Detected by] Me, when reading the final figures for this report.
  \item[Proof] In \texttt{report/figures/comparison\_overview.png} of 24 November, before the fix, the CLIP
        dot and bar lay under the legend box.
  \item[Correction] Legend outside the axes (\texttt{loc="center left", bbox\_to\_anchor=(1.01, 0.5)}),
        commit \texttt{032cea4}; figure regenerated and inspected; no number changed.
  \item[Reproduction] Put back \texttt{loc="lower left"} and run \texttt{python -m src.comparison report}.
  \item[Reflection] ``Figure inspected'' had meant checking that the numbers were right, not that every
        point was visible; a figure is checked as a reader sees it.
\end{description}

\subsubsection*{Sheet C3. Commands without the repository folder (Claude)}
\begin{description}[style=nextline,leftmargin=1.2em]
  \item[Prompt] What to run after the code of the learning-strategy block was delivered.
  \item[Faulty response] Activation of the environment and \texttt{python -m pytest -q tests} without
        \texttt{cd} into \texttt{ptbxl-dl}, although the terminal prompt showed the parent folder.
  \item[Detected by] Me, when the first command failed in my terminal.
  \item[Proof] PowerShell answered that \texttt{.venv\textbackslash Scripts\textbackslash activate} was
        not found and ``ERROR: file or directory not found: tests''.
  \item[Correction] Every later list of commands starts with \texttt{cd C:\textbackslash Users\textbackslash ori\textbackslash DL\textbackslash ptbxl-dl}.
  \item[Reproduction] From the parent folder, \texttt{python -m pytest -q tests} gives the same error.
  \item[Reflection] A small error, but it shows that the assistant reasons about the repository, not
        about the terminal the user is in.
\end{description}

\subsubsection*{Sheet C4. A wrong cause stated as fact (Claude)}
\begin{description}[style=nextline,leftmargin=1.2em]
  \item[Prompt] Why does the log of the trivial run say \texttt{git\_dirty: true}?
  \item[Faulty response] ``The run was made before the code was committed'', written as the cause in a
        block note.
  \item[Detected by] The assistant, from the \texttt{git status --short} output that I pasted from my
        Dense runs.
  \item[Proof] Every Dense run, all made after the commit with identical code, also said
        \texttt{git\_dirty: true}. From the \texttt{git status --short} output I pasted, the real cause
        appeared: each run writes its own predictions file before running \texttt{git status}, which lists
        untracked files, so every full run looked dirty.
  \item[Correction] The flag now looks only at the code under \texttt{src/} (log schema 2), with tests on a
        temporary repository; the old logs keep a note that their flag carries no information.
  \item[Reproduction] With a clean tree and the old check, \texttt{python -m src.train --model trivial
        --max-epochs 1 --tag check} logs \texttt{git\_dirty: true}.
  \item[Reflection] A plausible explanation was accepted without a second observation; the second block
        provided it.
\end{description}

\subsubsection*{Sheet C5. The smoke mode of the Dense report stopped (Claude)}
\begin{description}[style=nextline,leftmargin=1.2em]
  \item[Prompt] Write the complete Dense block, with an \texttt{all --fast} smoke test of the whole
        pipeline.
  \item[Faulty response] In \texttt{--fast} mode \texttt{build\_summary} paired the final runs with the
        trivial model on the test fold, but \texttt{--fast} runs see only 256 random test records and the
        trivial run the whole fold, so the paired bootstrap stopped (``differ in test\_ecg\_id''). The
        assistant's tests used a temporary run list without the trivial run, so they could not see it.
  \item[Detected by] The assistant, when no \texttt{\_fast} summary, tables or figures appeared among the
        outputs of my smoke run; it reproduced the error on a synthetic cache.
  \item[Proof] The same call on a synthetic cache with a full-fold trivial run raises
        \texttt{ValueError: ... differ in test\_ecg\_id}.
  \item[Correction] No trivial comparison in \texttt{--fast} mode (full runs use the same test records and
        were not affected); the test fixture lists a full-fold trivial run, plus a regression test.
  \item[Reproduction] Remove \texttt{and not fast} from the trivial-comparison condition of
        \file{src/dense_experiments.py::build_summary} and run the Dense tests
        (\file{python -m pytest -q tests/test_dense.py}); the test
        \file{test_fast_summary_skips_the_trivial_comparison} fails.
  \item[Reflection] A smoke test is only useful if its outputs are checked: a missing file was the
        symptom, not an error message.
\end{description}

\subsection*{Reflection}

Most errors of the assistant were not in the deep-learning code but at its edges: the device, the memory,
the folder, the copy of a file. They were found because we ran everything ourselves on the real hardware
and read the whole output. The rule that the assistant proposes and explains while we decide and execute
kept every result one that we can reproduce and explain.
